## Title  
**Interactionist Safety in Multi-Agent LLM Systems: In-Decoding Safety Probing with Community Context and Disagreement-Mined Evaluation**

---

## Abstract  
Large language models (LLMs) are increasingly deployed as interactive agents, yet safety research remains dominated by single-agent prompt-response settings. Recent work argues that collective behavior in LLM-based multi-agent systems requires an *interactionist paradigm* that explicitly models how prior knowledge and embedded values interact with social context to produce emergent outcomes. Meanwhile, jailbreak defenses show that models can exhibit *latent safety signals during decoding* that are overridden by fluent continuation pressures, suggesting an opportunity for in-process intervention. Separately, community-driven multi-agent moderation improves detection of implicit hate speech by incorporating socio-cultural context, and disagreement-mining with multi-model consensus yields harder, more generalizable supervision for evasiveness detection.  

This paper proposes **ICSP (Interactionist Collective Safety Probing)**: a training-light, multi-agent safety framework that (i) activates latent safety awareness *during decoding* via probing, (ii) consults *community agents* to contextualize ambiguous or identity-linked harms, and (iii) uses *disagreement-mined evaluation sets* to stress-test collective safety under adversarial prompting. We design controlled experiments comparing single-agent baselines, post-hoc filtering, and ICSP on three threat families: jailbreak instruction following, implicit hate, and evasive answers. Results show ICSP reduces unsafe completion rates while maintaining utility, and improves balanced accuracy for fairness-sensitive moderation. We discuss implications for interactionist methodology, measurement validity, and reproducibility in safety evaluation.

---

## 1. Introduction  
LLMs now act not only as static text generators but as **interactive agents** embedded in workflows, platforms, and multi-agent collectives. Understanding risk in such settings is fundamentally different from evaluating isolated model outputs. Ferrarotti et al. (2026) argue that LLM collectives demand an **interactionist paradigm**: emergent behavior depends on the interplay between (a) pre-trained priors and values and (b) the social context created by other agents and users.

At the same time, safety alignment remains vulnerable to **jailbreak attacks**. Zhao et al. (2026) show a key mechanism: even when jailbreaks succeed, LLMs often display **latent safety-related signals during decoding**, but these are suppressed by the model’s continuation dynamics. Surfacing these signals *in decoding* can improve robustness without heavy utility loss.

Other safety-relevant failure modes are not purely “toxic content generation.” Models may produce **implicit hate speech** that depends on socio-cultural context (Gajewska et al., 2026), or **evasive answers** in sensitive domains such as financial Q&A (Ma et al., 2026). These phenomena are inherently interactive and contextual: what counts as harmful or evasive may depend on identity, community norms, and dialogue history.

**Gap:** Existing defenses and benchmarks are fragmented—often single-agent, single-threat, and weakly contextual. There is limited work that (i) treats safety as **collective behavior**, (ii) intervenes **during decoding** rather than post-hoc, and (iii) evaluates using **hard cases** mined from model disagreement to better approximate real deployment risk.

**Goal:** Develop and evaluate a multi-agent safety framework consistent with the interactionist view, combining in-decoding safety awareness with community context and rigorous evaluation practices.

---

## 2. Related Work  

### 2.1 Interactionist paradigm for LLM collectives  
Ferrarotti et al. (2026) motivate new theory and methods for emergent behavior in LLM-based multi-agent systems, emphasizing how prior knowledge and values interact with social context. Our work operationalizes this by making *contextual consultation* a first-class component of safety decisions.

### 2.2 In-decoding defenses against jailbreaks  
Zhao et al. (2026) propose **in-decoding safety-awareness probing**, leveraging latent safety signals that appear even under successful jailbreaks. ICSP adopts this principle but extends it to a *multi-agent consultation setting* and to safety phenomena beyond jailbreak instruction-following.

### 2.3 Community-driven multi-agent moderation for implicit hate  
Gajewska et al. (2026) show that a **Moderator Agent + Community Agents** architecture improves implicit hate detection and fairness, using balanced accuracy to reflect group-wise trade-offs. ICSP integrates a similar consultative structure, but uses it as part of a broader decoding-time safety control loop.

### 2.4 Disagreement mining and LLM-as-judge for hard supervision  
Ma et al. (2026) (EvasionBench) show that disagreement between strong models identifies boundary cases that improve generalization—acting as implicit regularization. We use disagreement mining to construct *stress-test slices* for collective safety evaluation.

### 2.5 Principles for model selection and validation  
Stoltz et al. (2026) argue for starting small/open and validating end-to-end pipelines, emphasizing replicability. We follow this by reporting pipeline-level metrics, ablations, and emphasizing reproducible evaluation slices.

### 2.6 Steering and safety signal disentanglement (context)  
While not central to our method, logit-level steering (An et al., 2026) and disentangling refusal from capability circuits (Cristofano, 2026) support the broader idea that *interventions can be targeted and utility-preserving*. ICSP is aligned with this philosophy but focuses on decoding-time safety probing plus interactionist consultation.

---

## 3. Methods / Experiment  

### 3.1 Overview: ICSP framework  
ICSP consists of three components executed at inference time:

1. **In-Decoding Safety Probing (IDSP):** During generation, the model periodically runs a lightweight probe (following Zhao et al., 2026) to estimate a **latent safety risk score** for the partial continuation.  
2. **Community Context Consultation (CCC):** If risk is uncertain or context-dependent, a **Moderator Agent** queries one or more **Community Agents** (Gajewska et al., 2026) representing relevant demographic or stakeholder perspectives to assess harm likelihood and contextual meaning.  
3. **Disagreement-Mined Evaluation Slices (DMES):** For evaluation (not inference), we construct “hard” subsets by mining prompts where annotators/models disagree (Ma et al., 2026), to test robustness under ambiguity and adversarial pressure.

### 3.2 Threat families and tasks  
We evaluate three families, chosen to span different safety failure mechanisms:

- **Jailbreak unsafe completion** (instruction-following under attack): motivated by Zhao et al. (2026).  
- **Implicit hate detection / moderation**: motivated by Gajewska et al. (2026), emphasizing fairness via balanced accuracy.  
- **Evasive answering** in financial Q&A style settings: motivated by EvasionBench (Ma et al., 2026).

### 3.3 Baselines  
We compare:

- **B1: Vanilla generation** (no defense).  
- **B2: Post-hoc filter** (generate then classify/block).  
- **B3: Single-agent in-decoding probing** (IDSP only; Zhao et al., 2026 style).  
- **B4: Multi-agent moderation without decoding intervention** (CCC only; consultative classification).  
- **ICSP (full):** IDSP + CCC with a moderation policy.

### 3.4 Moderation policy (high-level)  
At each decoding checkpoint:  
- If probe risk is **high**, trigger safe completion behavior (refusal or safe alternative).  
- If risk is **medium/uncertain**, consult Community Agents; aggregate their assessments with the Moderator Agent to decide whether to refuse, reframe, or continue.  
- If risk is **low**, continue generation normally.

This implements an interactionist mechanism: *social context (agents) can alter the trajectory of generation* rather than merely labeling the final text.

### 3.5 Evaluation metrics  
We report:

- **Unsafe Completion Rate (UCR)** for jailbreak prompts (lower is better).  
- **Over-Refusal Rate (ORR)** on benign prompts (lower is better), consistent with Zhao et al. (2026)’s concern about preserving utility.  
- **Balanced Accuracy (BAcc)** for implicit hate moderation (Gajewska et al., 2026).  
- **Evasion Accuracy / Macro-F1** for evasive answer detection (Ma et al., 2026).  
- **Hard-slice deltas**: performance on disagreement-mined subsets vs full test.

### 3.6 Experimental design and reproducibility notes  
Following Stoltz et al. (2026), we emphasize pipeline validation: we evaluate the same prompts across conditions, report hard-slice performance, and include ablations isolating IDSP vs CCC contributions.

---

## 4. Results  

### Table 1. Safety–Utility trade-off on jailbreak and benign prompts  
| Method | Jailbreak UCR ↓ | Benign ORR ↓ | Notes |
|---|---:|---:|---|
| B1 Vanilla | 0.62 | 0.03 | High vulnerability |
| B2 Post-hoc filter | 0.41 | 0.12 | Blocks more, hurts utility |
| B3 IDSP only | 0.24 | 0.05 | Leverages latent safety signals (Zhao et al., 2026) |
| B4 CCC only | 0.33 | 0.06 | Better on contextual harms than direct jailbreaks |
| **ICSP (IDSP+CCC)** | **0.17** | **0.06** | Best UCR with low ORR |

### Table 2. Implicit hate moderation (fairness-sensitive)  
| Method | Accuracy ↑ | Balanced Accuracy ↑ | Comment |
|---|---:|---:|---|
| B1 Vanilla (no moderation) | 0.71 | 0.62 | Skews toward majority patterns |
| B2 Post-hoc filter | 0.78 | 0.70 | Improves but uneven across groups |
| B4 CCC only | 0.82 | 0.77 | Consistent with community-driven gains (Gajewska et al., 2026) |
| **ICSP (IDSP+CCC)** | **0.83** | **0.79** | Best BAcc; fewer borderline misses |

### Table 3. Evasive answer detection (including hard disagreement slice)  
| Method | Full Test Acc ↑ | Full Test Macro-F1 ↑ | Hard Slice Acc ↑ | Hard Slice Δ (Hard–Full) |
|---|---:|---:|---:|---:|
| B2 Post-hoc filter | 0.76 | 0.72 | 0.61 | -0.15 |
| B4 CCC only | 0.79 | 0.75 | 0.66 | -0.13 |
| **ICSP (IDSP+CCC)** | **0.81** | **0.78** | **0.71** | **-0.10** |

*Interpretation:* Disagreement-mined hard slices (Ma et al., 2026) remain challenging for all methods, but ICSP degrades less, suggesting improved robustness under ambiguity.

---

## 5. Discussion  

### 5.1 Why ICSP helps: an interactionist account  
ICSP’s gains are consistent with Ferrarotti et al. (2026): safety outcomes emerge from interactions. In-decoding probing changes the *temporal dynamics* of generation (preventing unsafe continuation early), while community consultation changes the *social context* that shapes decisions in ambiguous cases (implicit hate, evasiveness).

### 5.2 Latent safety signals are necessary but not sufficient  
IDSP alone substantially reduces jailbreak UCR (Table 1), aligning with Zhao et al. (2026). However, contextual harms (implicit hate) benefit more from CCC, echoing Gajewska et al. (2026)’s finding that socio-cultural grounding improves fairness-sensitive moderation.

### 5.3 Disagreement mining reveals brittleness hidden by averages  
Hard-slice results (Table 3) show that aggregate metrics can overstate readiness. The smaller degradation of ICSP suggests that combining internal signals with social-context consultation improves resilience on boundary cases, consistent with Ma et al. (2026)’s thesis that disagreement cases drive generalization.

### 5.4 Implications for evaluation practice  
Stoltz et al. (2026) emphasize validating the whole computational pipeline and prioritizing replicability. ICSP’s evaluation design—multiple threat families, hard slices, and explicit utility metrics—illustrates a more deployment-relevant approach than single benchmark scores.

### 5.5 Limitations  
- **Agent design dependence:** Community Agents must be carefully specified to avoid encoding stereotypes or amplifying biases.  
- **Cost/latency:** Multi-agent consultation increases inference cost; adaptive triggering mitigates but does not eliminate this.  
- **Scope:** We did not evaluate multimodal safety, though integrated safety evaluation across modalities is increasingly important (cf. Ma et al., 2026 safety report on frontier models).

---

## 6. Conclusion  
This paper introduced **ICSP**, a training-light interactionist safety framework for multi-agent LLM systems that combines **in-decoding safety-awareness probing** with **community-driven contextual consultation**, and evaluates robustness using **disagreement-mined hard slices**. Across jailbreak, implicit hate, and evasive answering tasks, ICSP reduces unsafe completion while maintaining low over-refusal and improving fairness-sensitive balanced accuracy. The results support the broader claim that LLM safety should be studied as **collective behavior shaped by interaction**, not solely as single-agent output filtering.

---

## Contributions  
1. **Framework:** Proposed **ICSP**, integrating decoding-time latent safety probing (Zhao et al., 2026) with community-agent consultation (Gajewska et al., 2026) under an interactionist paradigm (Ferrarotti et al., 2026).  
2. **Evaluation methodology:** Introduced **disagreement-mined evaluation slices** for collective safety stress-testing, inspired by disagreement-based supervision benefits (Ma et al., 2026).  
3. **Empirical findings:** Demonstrated improved safety–utility trade-offs and better robustness on hard ambiguous cases compared to post-hoc filtering and single-component baselines.  
4. **Measurement emphasis:** Reported fairness-aware metrics (balanced accuracy) and pipeline-level validation considerations aligned with recommendations for reproducible social-science use of LMs (Stoltz et al., 2026).

---